\documentclass[11pt]{article}

\usepackage[a4paper,margin=1in]{geometry}
\usepackage{amsmath,amssymb,amsthm,mathtools}
\usepackage{booktabs}
\usepackage{enumitem}
\usepackage{microtype}
\usepackage{xcolor}
\usepackage{hyperref}

\hypersetup{
    colorlinks=true,
    linkcolor=blue!60!black,
    citecolor=blue!60!black,
    urlcolor=blue!60!black,
    pdftitle={A Golden Shadow Algebra},
    pdfauthor={Vinicius Santos}
}

\newtheorem{definition}{Definition}[section]
\newtheorem{theorem}[definition]{Theorem}
\newtheorem{corollary}[definition]{Corollary}
\newtheorem{proposition}[definition]{Proposition}
\newtheorem{lemma}[definition]{Lemma}
\newtheorem{observation}[definition]{Observation}
\newtheorem{conjecture}[definition]{Conjecture}
\newtheorem{question}[definition]{Question}
\newtheorem{remark}[definition]{Remark}

\newcommand{\B}{\mathsf{B}}
\newcommand{\Bk}{\mathsf{B}_k}
\newcommand{\Scl}{\mathsf{Scal}}
\newcommand{\Pair}{\mathsf{Pair}}
\newcommand{\Tr}{\operatorname{Tr}}
\newcommand{\Nm}{\operatorname{N}}
\newcommand{\Sp}{\operatorname{Sp}}
\newcommand{\sig}{\sigma}
\newcommand{\ph}{\varphi}
\newcommand{\ps}{\psi}
\newcommand{\Q}{\mathbb{Q}}
\newcommand{\Z}{\mathbb{Z}}
\newcommand{\R}{\mathbb{R}}

\title{A Golden Shadow Algebra\\
\large From a $\varphi$-Leaf Scalar Core to Conjugate-Pair Arithmetic}
\author{Vinicius Santos\\
\small Exploratory draft}
\date{June 2026}

\begin{document}
\maketitle

\begin{abstract}
The recent EML construction of Odrzywo\l{}ek shows that the elementary functions of a scientific calculator can be generated from a single binary operator,
\[
\operatorname{eml}(x,y)=e^x-\log y,
\]
together with the distinguished leaf $1$.  This note asks a complementary algebraic question: what kind of generated structure arises if the distinguished leaf is not the rational unit $1$, but the golden ratio
\[
\ph=\frac{1+\sqrt5}{2}?
\]
Guided by exact symbolic searches, we isolate the scalar operation
\[
\B(x,y)=xy-y=y(x-1),
\]
which contains no occurrence of $\ph$ in its definition.  With leaf $\ph$, this operation recovers the unit through $\B(\ph,\ph)=1$ and generates a scalar core of predecessor, negation, successor, golden contraction, and multiplication.  However, binary addition is not merely hard to find: it is provably absent from the scalar grammar.  More precisely, the scalar fragment is exactly the multiplicative closure of one-variable polynomials under constant shifts from $\Z[\ph]$.

The obstruction suggests that the missing operation belongs to another sort.  We introduce a two-sorted shadow algebra with scalar terms and pair terms $\langle x\mid y\rangle$.  The primitive pair-level additive channel is the trace projection
\[
\Tr\langle x\mid y\rangle=x+y.
\]
The main content is therefore a scalar characterization and an incompleteness/completion pair: the scalar fragment cannot express genuine mixed addition, while adjoining pair formation and trace gives a natural two-sorted completion of ordinary ring arithmetic; multiplication, norm, and spread are derived operations.  The exact scalar constants are $\Z[\ph]$: they are countable and therefore not all real numbers, but under a chosen real embedding they are dense in $\R$, so their metric completion is the real line.  Under the full Minkowski embedding into $\R^2$, the same ring is discrete.  Thus the construction is exact-algebraic at finite depth and real-analytic only after completion.
\end{abstract}

\tableofcontents

\section{Introduction: what does a leaf know?}

A striking feature of the EML construction is its economy.  Odrzywo\l{}ek~\cite{Odrzywolek2026EML} introduced the binary operation
\[
    \operatorname{eml}(x,y)=e^x-\log y
\]
and showed, constructively, that together with the single constant $1$ it generates the standard repertoire of a scientific calculator.  In this setting, every expression becomes a binary tree with one node type:
\[
    S \to 1 \mid \operatorname{eml}(S,S).
\]

The present note begins with a simple variation on this theme.  What if the distinguished leaf is not the rational unit $1$, but the golden ratio
\[
    \ph=\frac{1+\sqrt5}{2}?
\]
The rational unit $1$ is operationally neutral.  The golden ratio is different: it is not neutral, but it is algebraically rich.  It satisfies
\[
    \ph^2-\ph=1,
\]
so the rational unit is not assumed; it is generated by $\ph$ interacting with itself.  The golden ratio is also extremal in Diophantine approximation through its continued fraction
\[
    \ph=[1;1,1,1,\dots],
\]
which explains its role as a worst-case irrational in Hurwitz-type approximation bounds~\cite{HardyWright2008,Hurwitz1891}.

The guiding question is therefore:
\begin{quote}
If EML explores a one-operator universe generated from the rational unit $1$, what kind of algebraic universe is naturally generated from the extremal irrational leaf $\ph$?
\end{quote}

The answer developed here is not a new proof of elementary-function universality.  Instead, the exploration points toward a different structure: a two-layer algebra in which a clean scalar operation generates a multiplicative/predecessor core, while ordinary addition appears only after introducing a conjugate-pair shadow layer.

The main results are as follows.  First, the scalar core generates basic arithmetic structure from the golden leaf, including $1$, $0$, negation, predecessor, successor, and multiplication.  Second, mixed affine addition is not scalar-definable.  Third, the scalar fragment is characterized exactly as the multiplicative, constant-shift closure of one-variable polynomial rings over $\Z[\ph]$.  Fourth, pair formation together with trace gives a natural two-sorted additive completion.

\section{The golden seed}

Let
\[
    \ph=\frac{1+\sqrt5}{2},\qquad \ps=\frac{1-\sqrt5}{2}.
\]
Then $\ph$ and $\ps$ are the two roots of
\[
    t^2-t-1=0.
\]
The nontrivial Galois conjugation of $\Q(\sqrt5)/\Q$ sends
\[
    \sig(\ph)=\ps=1-\ph=-\ph^{-1}.
\]
The basic trace and norm identities are
\[
    \ph+\ps=1,
\]
\[
    \ph\ps=-1,
\]
and the spread between the two embeddings is
\[
    \ph-\ps=2\ph-1.
\]
Under the chosen positive real embedding, this spread is the positive square root $\sqrt5$.

These facts are standard algebraic number theory: the ring of integers of $\Q(\sqrt5)$ is $\Z[\ph]$, and conjugation, trace, and norm are the canonical maps associated to this quadratic extension~\cite{Marcus1977,ConradTraceNorm}.  What is less standard is to reinterpret them as operations in a generated-expression system.

\section{Why naive \texorpdfstring{$\varphi$}{phi}-operators became complicated}

The first attempts at a $\ph$-leaf analogue of EML naturally imitated the analytic shape of EML\@.  One can write down shifted $\exp/\log$ variants, hyperbolic variants involving $\tanh/\operatorname{atanh}$, or branch variants involving $\tan/\arctan$.  Some of them are genuinely suggestive.  For instance,
\[
    A_\ph(x,y)=\frac{x}{\ph}-\arctan(\tan(y-\ph))
\]
quickly exposes a local affine chart, because on a principal branch it behaves like
\[
    A_\ph(x,y)=\frac{x}{\ph}-y+\ph,
\]
and it can even expose $\pi$ through branch jumps.  But this is already suspicious: the formula contains $\ph$ in the operator, imports a branch mechanism, and is visibly tuned to make ordinary addition appear.

That discomfort became the first lesson of the exploration.  A true leaf-based analogue should not smuggle the leaf into the operation.  EML is not written as $\operatorname{eml}_1$; the operator is independent of the constant $1$.  We therefore impose the same discipline here:
\begin{quote}
The operator should be independent of $\ph$; only the leaf should be $\ph$.
\end{quote}
The second lesson is just as important.  Earlier formulas became complicated because a single scalar expression was being asked to do two jobs at once: generate multiplication-like structure and generate mixed addition.  The clean scalar reset below lets the first job happen naturally, and then proves that the second job belongs somewhere else.

\section{The scalar reset: predecessor-scaling}

\begin{definition}[The scalar core]
Define the binary scalar operation
\[
    \B(x,y)=xy-y=y(x-1).
\]
The scalar grammar generated by $\B$ and the leaf $\ph$ is
\[
    S\to \ph \mid x \mid y \mid \cdots \mid \B(S,S),
\]
where variables are included when studying definable operations.
\end{definition}

\begin{proposition}[Basic scalar constructions]\label{prop:scalar-constructions}
In the scalar grammar generated by $\B$ and $\ph$, the following terms are definable:
\[
    1,\quad 0,\quad \ph^{-1},\quad \ps,\quad -x,\quad x-1,\quad x+1,\quad xy.
\]
\end{proposition}

\begin{proof}
The unit is obtained from the defining identity of $\ph$:
\[
    1=\B(\ph,\ph)=\ph^2-\ph.
\]
Then
\[
    0=\B(1,\ph)=\ph-\ph.
\]
The inverse golden unit is
\[
    \ph^{-1}=\B(\ph,1)=\ph-1.
\]
The conjugate is
\[
    \ps=-\ph^{-1}=\B(0,\B(\ph,1)).
\]
For a variable $x$,
\[
    -x=\B(0,x),
\]
and
\[
    x-1=\B(x,1).
\]
A successor term is obtained by negating the predecessor of $-x$:
\[
    x+1=\B\bigl(0,\B(\B(0,x),1)\bigr).
\]
Finally,
\[
    xy=\B(x+1,y),
\]
since $\B(x+1,y)=y((x+1)-1)=xy$.
\end{proof}

The operation has a simple native interpretation:
\[
    \B(x,y)=\text{``scale }y\text{ by the predecessor of }x\text{''}.
\]
In particular,
\[
    \B(\ph,y)=y(\ph-1)=\frac{y}{\ph}.
\]
Thus $\ph$ acts by golden contraction in the visible scalar layer.

\section{The scalar fragment}

The same scalar grammar that produces multiplication naturally resists binary addition.  This failure is not merely computational.  It is forced by the factorization pattern of the scalar operation.

\begin{theorem}[Scalar non-additivity]\label{thm:scalar-nonadditivity}
Let $K$ be a field and let $\mathcal T_B$ be the scalar term grammar generated from constants in $K$, variables $x,y$, and the single binary operation
\[
    B(f,g)=fg-g=g(f-1).
\]
Then no term of $\mathcal T_B$ denotes a mixed affine-linear polynomial
\[
    ax+by+c
\]
with $a,b\in K^\times$.  In particular, neither $x+y$ nor $x-y$ is generated by the scalar grammar.
\end{theorem}

\begin{proof}
Assume that such a term exists, and choose one of minimal term size.  It cannot be a leaf, since the leaves are constants, $x$, and $y$, and none of these contains both variables linearly.  Thus the term has the form
\[
    B(f,g)=g(f-1).
\]
The target $ax+by+c$ is nonzero and has total degree $1$.  Over a field, total degree is additive on nonzero products:
\[
    \deg(uv)=\deg u+\deg v.
\]
Therefore one of the factors $g$ or $f-1$ must be constant.  If $g$ is a nonzero constant, then $f-1$ is again a mixed affine-linear polynomial, and hence $f$ is a smaller term denoting a mixed affine-linear polynomial.  If $f-1$ is a nonzero constant, then $g$ is a smaller term denoting a mixed affine-linear polynomial.  Both cases contradict minimality.
\end{proof}

The same proof applies to the shifted cores
\[
    B_k(f,g)=fg-kg=g(f-k)
\]
for any fixed $k\in K$.

\begin{corollary}[Non-additivity for shifted scalar cores]\label{cor:shifted-nonadditivity}
Let $K$ be a field and fix $k\in K$.  In the scalar grammar generated by constants, variables, and
\[
    B_k(f,g)=g(f-k),
\]
no mixed affine-linear polynomial $ax+by+c$ with $a,b\in K^\times$ is definable.
\end{corollary}

\subsection{The scalar fragment characterized}\label{subsec:scalar-characterization}

The non-additivity theorem rules out the most visible missing operation, but it does not yet say what the scalar fragment can express.  The answer is clean: scalar terms are exactly the polynomials obtained from one-variable pieces by multiplying and shifting by constants.

Let
\[
    R=\Z[\ph].
\]
Define $\mathcal M\subset R[x,y]$ to be the smallest collection of polynomials satisfying:
\begin{enumerate}[label=(\roman*)]
    \item $R[x]\subset \mathcal M$ and $R[y]\subset \mathcal M$;
    \item if $P,Q\in\mathcal M$, then $PQ\in\mathcal M$;
    \item if $P\in\mathcal M$ and $c\in R$, then $P+c\in\mathcal M$.
\end{enumerate}
Thus $\mathcal M$ is the multiplicative, constant-shift closure of the one-variable polynomial rings.  Mixed dependence can arise through products such as $p(x)q(y)$, but not through direct additive mixing such as $p(x)+q(y)$.

\begin{theorem}[Scalar fragment characterization]\label{thm:scalar-characterization}
Let $D$ be the set of polynomials in $R[x,y]$ definable by the scalar grammar generated from the leaf $\ph$, variables $x,y$, and
\[
    \B(f,g)=g(f-1).
\]
Then
\[
    D=\mathcal M.
\]
Moreover, the scalar-definable constants are exactly $R=\Z[\ph]$.
\end{theorem}

\begin{proof}
We separate the proof into three steps.

\emph{Step 1: scalar constants.}  Every scalar term has coefficients in $R$, since the leaf $\ph$ lies in $R$ and $\B(f,g)=fg-g$ uses only multiplication and subtraction with integer coefficients.  Hence scalar-definable constants are contained in $R$.

The constructions of Proposition~\ref{prop:scalar-constructions} give $1$, $0$, negation, predecessor, successor, and multiplication.  In particular, for every scalar-definable $P$, the terms $P+1$, $P-1$, and $PQ$ are scalar-definable.  Since
\[
    \ph^{-1}=\ph-1
\]
was also constructed, we can add $\ph$ to any scalar-definable $P$ by
\[
    P+\ph=\B(P\ph^{-1}+2,\ph).
\]
Repeatedly applying integer shifts and this $\ph$-shift shows that $P+c$ is scalar-definable for every $c\in R$.  Taking $P=0$ shows that every element of $R$ is a scalar-definable constant.  Thus the scalar-definable constants are exactly $R$.

\emph{Step 2: closure and the inclusion $D\subseteq\mathcal M$.}  The inclusion follows by induction on terms: the leaves $\ph,x,y$ belong to $\mathcal M$, and if $f,g\in\mathcal M$, then
\[
    \B(f,g)=g(f-1)
\]
belongs to $\mathcal M$ because $f-1\in\mathcal M$ and $\mathcal M$ is closed under multiplication and constant shifts.

\emph{Step 3: one-variable generation and the inclusion $\mathcal M\subseteq D$.}  Scalar terms are closed under multiplication and under addition of constants from $R$.  Every one-variable polynomial in $R[x]$ is scalar-definable by Horner's rule: starting from the leading coefficient, repeatedly multiply by $x$ and add the next coefficient.  The same argument gives all of $R[y]$.  Since $D$ contains the generators of $\mathcal M$ and is closed under the operations defining $\mathcal M$, we have $\mathcal M\subseteq D$.
\end{proof}

Before adding the shadow layer, we already know exactly what the visible scalar layer can do: it builds multiplicative trees whose leaves are one-variable polynomials, and it allows constant shifts at any node.  Additive mixing of independent variables can appear only as a byproduct of multiplication, never as a primitive scalar operation.

\subsection{Exact generation versus completion}\label{subsec:completion}

The scalar characterization also clarifies a natural question: can the mechanism generate all real numbers?

There are two different meanings of ``generate.''  If generation means finite scalar terms, then the answer is no.  For constants, the exact finite scalar constants are
\[
    \Z[\ph]=\{a+b\ph:a,b\in\Z\},
\]
by Theorem~\ref{thm:scalar-characterization}.  This set is countable, whereas $\R$ is uncountable.  Therefore no finite expression-tree grammar of this kind can generate every real number exactly.

However, after choosing the embedding $\ph=(1+\sqrt5)/2\in\R$, the exact constants are dense in the real line.  This is a one-embedding statement: under the full Minkowski embedding $a+b\ph\mapsto(a+b\ph,a+b\ps)$ into $\R^2$, the same ring $\Z[\ph]$ is a discrete lattice.  The density used here is topological rather than computationally cheap: close approximants may require large coefficients, and finitely described limiting procedures still form only a countable class.

\begin{proposition}[Dense golden constants]\label{prop:dense-golden-constants}
Under the chosen real embedding $\ph=(1+\sqrt5)/2$, the scalar-definable constants do not equal $\R$, but their closure in the usual topology is all of $\R$:
\[
    \overline{\Z[\ph]}=\R.
\]
\end{proposition}

\begin{proof}
The scalar-definable constants are exactly $\Z[\ph]$, hence are countable and cannot equal $\R$.

To prove density, let $r\in\R$ and $\varepsilon>0$.  Since $\ph$ is irrational, the fractional parts of $b\ph$, with $b\in\Z$, are dense modulo $1$~\cite{HardyWright2008}.  Hence we can choose $b\in\Z$ such that $b\ph$ is within $\varepsilon$ of $r$ modulo $1$.  Choosing $a\in\Z$ to correct the integer part gives
\[
    |a+b\ph-r|<\varepsilon.
\]
Thus every real number is a limit of scalar-definable constants in this selected embedding.  This does not contradict the discreteness of $\Z[\ph]$ under the full Minkowski embedding into $\R^2$.
\end{proof}

So the mechanism has three levels:
\[
    \text{finite scalar constants}=\Z[\ph],
\]
\[
    \text{topological closure}=\R,
\]
and
\[
    \text{computably described limits}\subsetneq\R.
\]
The last distinction matters: allowing arbitrary Cauchy limits gives all real numbers, while allowing only finitely describable or computable limiting processes still gives only countably many individually describable reals.  Thus the correct statement is not that the grammar generates every real number as a finite expression.  Rather, in a chosen real embedding it generates a dense golden integer ring whose metric completion is the real line.

\subsection{Descent and irreducibility consequences}

The characterization gives a recursive test.  If $P\in R[x,y]$ depends on both variables, then $P\in D$ precisely when either $P$ is already a one-variable polynomial, or there exists $c\in R$ such that
\[
    P-c=Q_1Q_2
\]
with $Q_1,Q_2\in D$.  Iterating this condition terminates by degree descent whenever nonconstant factors occur.  Constant factors simply record multiplication by definable constants in $R$.

Theorem~\ref{thm:scalar-nonadditivity} is the first-degree case of this descent.  The next obstruction handles many higher-degree additive-looking targets.

\begin{proposition}[Affine-shift irreducibility obstruction]\label{prop:irreducibility-obstruction}
Let $P\in R[x,y]$.  Suppose that for every $c\in R$, the polynomial $P-c$ is irreducible in $R[x,y]$ and depends nontrivially on both variables.  Then $P$ is not scalar-definable.
\end{proposition}

\begin{proof}
If $P\in D=\mathcal M$ and $P$ depends on both variables, the descent criterion gives a constant $c\in R$ such that $P-c$ factors as a product of scalar-definable terms.  This contradicts the assumed irreducibility of every constant shift $P-c$.
\end{proof}

\begin{corollary}[Some separable sums are scalar-impossible]\label{cor:separable-sums}
The polynomials
\[
    x^m+y\qquad\text{and}\qquad x+y^m
\]
are not scalar-definable for any $m\ge 1$.  The polynomial $x^2+y^2$ is also not scalar-definable over $R=\Z[\ph]$.
\end{corollary}

\begin{proof}
For $P=x^m+y$, every constant shift $P-c=x^m+y-c$ is linear in $y$ with unit leading coefficient.  Viewed as a polynomial in $y$ over $R[x]$, any factorization would force a nonconstant factor of $R[x]$ to divide the unit coefficient of $y$, which is impossible.  The same argument applies to $x+y^m$.

For $P=x^2+y^2$, if $x^2+y^2-c$ factored over $R[x,y]$, then over the fraction field $\Q(\ph)(x)$ it would factor as a quadratic in $y$.  This would require $-(x^2-c)$ to be a square in $\Q(\ph)(x)$.  Comparing leading terms shows that $-1$ would have to be a square in $\Q(\ph)$, impossible because $\Q(\ph)\subset\R$.
\end{proof}

Irreducibility is not the whole story.  Reducible polynomials can still fail the descent criterion.  For example,
\[
    x^2-y^2=(x-y)(x+y)
\]
is reducible, but its nontrivial factors are the mixed affine-linear polynomials $x-y$ and $x+y$, both excluded by Theorem~\ref{thm:scalar-nonadditivity}.  Constant shifts do not rescue the descent.  If $c\ne0$, then $x^2-y^2-c$ is irreducible over $R[x,y]$: over the fraction field and after the linear change of variables $u=x-y$, $v=x+y$, it becomes $uv-c$, which is irreducible.  If $c=0$, the only useful linear factors are precisely $x-y$ and $x+y$, which fail the scalar criterion.  Thus $x^2-y^2$ is scalar-impossible even though it is reducible.

Rather than forcing addition into the scalar core, we next ask whether addition belongs to another layer.

\section{The conjugate shadow}

The golden ratio is not alone.  It has a conjugate shadow
\[
    \ps=1-\ph=-\ph^{-1}.
\]
A general element of $\Q(\ph)$ can be written as $a+b\ph$, and its conjugate is
\[
    (a+b\ph)^\sig=a+b\ps.
\]
Thus a $\ph$-number naturally has two embeddings:
\[
    x\longmapsto (x,x^\sig).
\]
The standard trace and norm maps are
\[
    \Tr(x)=x+x^\sig,
\]
\[
    \Nm(x)=xx^\sig.
\]
For $\ph$ itself,
\[
    \Tr(\ph)=\ph+\ps=1,
\]
\[
    \Nm(\ph)=\ph\ps=-1.
\]
These identities suggest that the rational unit $1$ can be obtained in two different ways:
\[
    1=\ph^2-\ph
\]
inside the visible scalar layer, and
\[
    1=\ph+\ps
\]
as a trace from the conjugate-pair layer.  This is the key conceptual shift: addition may not be a visible scalar operation in the $\ph$-core; it may be a projection from a conjugate pair.

\section{The typed shadow algebra}

The previous obstruction says that addition cannot be recovered by looking harder inside the scalar grammar.  The next move is therefore not to complicate $\B$, but to add a natural new kind of object that can merge two scalar expressions additively.

\begin{definition}[Golden shadow algebra]
The golden shadow algebra has two sorts:
\begin{enumerate}[label=(\roman*)]
    \item scalar terms $s,t\in\Scl$;
    \item pair terms $p,q\in\Pair$.
\end{enumerate}
Its scalar data are the leaf $\ph$ and the operation
\[
    \B(s,t)=st-t.
\]
Its pair-level primitive data are pair formation
\[
    \langle s\mid t\rangle\in\Pair
\]
and the trace projection
\[
    \Tr\langle s\mid t\rangle=s+t.
\]
We also record the pair involution
\[
    \sig\langle s\mid t\rangle=\langle t\mid s\rangle
\]
as the formal shadow symmetry.  It is conceptually important, but it is not needed to complete ordinary ring arithmetic.
\end{definition}

The point is not that the displayed formula for $\Tr$ is mysterious.  It is a primitive projection.  The content is the following incompleteness/completion pair: the scalar fragment provably cannot form mixed addition, while adjoining pair formation and one projection closes exactly that gap.

\begin{definition}[Derived scalar macros]\label{def:derived-macros}
Inside the scalar fragment, define
\[
    1_\ph=\B(\ph,\ph),\qquad 0_\ph=\B(1_\ph,\ph),
\]
\[
    \operatorname{neg}(x)=\B(0_\ph,x),\qquad \operatorname{pred}(x)=\B(x,1_\ph),
\]
\[
    \operatorname{succ}(x)=\operatorname{neg}(\operatorname{pred}(\operatorname{neg}(x))),
\]
and
\[
    \operatorname{mul}(x,y)=\B(\operatorname{succ}(x),y).
\]
By Proposition~\ref{prop:scalar-constructions}, these macros denote $1$, $0$, $-x$, $x-1$, $x+1$, and $xy$, respectively.
\end{definition}

\begin{theorem}[Scalar incompleteness and typed completion]\label{thm:completion}
The scalar fragment generated by $\{\B,\ph\}$ cannot define mixed affine-linear operations such as $x+y$ or $x-y$.  After adjoining pair formation and the trace projection, the typed system defines ordinary ring arithmetic:
\[
    0,\quad 1,\quad -x,\quad x+y,\quad x-y,\quad xy.
\]
Moreover, multiplication is already scalar-definable; the pair layer contributes exactly the missing additive projection.
\end{theorem}

\begin{proof}
The incompleteness statement is Theorem~\ref{thm:scalar-nonadditivity}.  The scalar constants and unary operations $0$, $1$, and $-x$ are given by Definition~\ref{def:derived-macros}.  Addition is the trace projection
\[
    x+y=\Tr\langle x\mid y\rangle.
\]
Subtraction is trace against scalar negation:
\[
    x-y=\Tr\langle x\mid \operatorname{neg}(y)\rangle.
\]
Finally, multiplication is already scalar-definable:
\[
    xy=\operatorname{mul}(x,y)=\B(\operatorname{succ}(x),y).
\]
Thus pair formation together with $\Tr$ completes the ordinary ring operations that the scalar fragment lacks.
\end{proof}

This theorem is the center of the paper.  The two-sorted layer is not a decorative way to rewrite $x+y$.  It gives a natural two-sorted completion of a scalar grammar that is provably non-additive.

\begin{definition}[Derived norm and spread]
For a pair $\langle s\mid t\rangle$, define the derived norm and spread by
\[
    \Nm\langle s\mid t\rangle=\operatorname{mul}(s,t),
\]
\[
    \Sp\langle s\mid t\rangle=\Tr\langle s\mid \operatorname{neg}(t)\rangle=s-t.
\]
When $t=s^\sig$, the derived norm agrees with the usual field norm.
\end{definition}

\begin{definition}[The golden pair]
Let
\[
    \Phi=\langle\ph\mid\ps\rangle.
\]
\end{definition}

\begin{proposition}[Golden projections]
The golden pair satisfies
\[
    \Tr(\Phi)=1,
\]
\[
    \Nm(\Phi)=-1,
\]
\[
    \Sp(\Phi)=\sqrt5.
\]
\end{proposition}

\begin{proof}
These are exactly the identities
\[
    \ph+\ps=1,
\]
\[
    \ph\ps=-1,
\]
\[
    \ph-\ps=\sqrt5.
\]
The norm identity uses the derived multiplication macro of Definition~\ref{def:derived-macros}.
\end{proof}

\begin{remark}[Vieta viewpoint]\label{rem:vieta}
A pair $\langle r_1\mid r_2\rangle$ determines a quadratic polynomial
\[
    t^2-\Tr\langle r_1\mid r_2\rangle t+\Nm\langle r_1\mid r_2\rangle.
\]
For $\Phi$, this polynomial is
\[
    t^2-t-1,
\]
the defining polynomial of the golden ratio.  This Vieta viewpoint is the bridge to the quadratic-unit generalization in Section~\ref{sec:quadratic-units}.
\end{remark}

\section{A small computational search}

The clean scalar core was discovered by exact symbolic search before the obstruction theorem was proved.  The search represented expressions as polynomials over $\Q(\ph)$ using
\[
    \ph^2=\ph+1,
\]
and compared reduced polynomials exactly rather than numerically.  Historically, the search served as a diagnostic: it found $1$, $0$, $\ph^{-1}$, $\ps$, $-x$, $x+1$, $2$, $\sqrt5$, and $xy$ at modest costs, while repeatedly failing to find $x+y$ or $x-y$.  Theorems~\ref{thm:scalar-nonadditivity} and~\ref{thm:scalar-characterization} now explain why: those targets are not merely deeper; they are outside the scalar fragment.

A parallel sweep over bilinear $\ph$-free operators
\[
    O(x,y)=a xy+b x+c y,\qquad a,b,c\in\{-1,0,1\},
\]
identified $xy-y$ and its mirror $xy-x$ as the strongest clean multiplicative/predecessor candidates.  Additive operators such as $x-y$ and $x+y$ generated addition-like structure immediately, but they did not expose the same golden scalar core.  The search therefore played the role of a compass: it pointed toward a simple scalar operation, and the proof later explained the boundary of that operation.

\section{Comparison with EML}

The EML result and the golden shadow algebra answer different questions.

\begin{center}
\begin{tabular}{lll}
\toprule
Feature & EML & Golden shadow algebra \\
\midrule
Leaf & $1$ & $\ph$ or $\Phi=\langle\ph\mid\ps\rangle$ \\
Main operation & $e^x-\log y$ & $\B(x,y)=xy-y$ \\
Sorts & one-sorted & two-sorted: scalar and pair \\
Main strength & elementary-function universality & algebraic $\ph$-native structure \\
Addition & reconstructed inside EML trees & projection $\Tr\langle x\mid y\rangle$ \\
Multiplication & reconstructed inside EML trees & scalar $\B$; $\Nm$ is derived \\
$\pi,e,\log,\sin$ & constructively generated & not generated here \\
Interpretation & analytic universal gate & conjugate-shadow arithmetic \\
\bottomrule
\end{tabular}
\end{center}

Thus the golden shadow algebra should not be presented as an improvement over EML in the original sense.  EML is a one-sorted analytic universality theorem.  The present construction is a two-sorted algebraic presentation.  Exact finite generation is countable in either setting; questions about the real line require density, completion, or limit semantics.

\section{What the formulation buys us}

The formulation is useful because it turns an apparent failure into a structural theorem.  The scalar operation
\[
    \B(x,y)=xy-y
\]
knows how to make $1$, $0$, negation, successor, and multiplication from the $\ph$ leaf, but Theorem~\ref{thm:scalar-nonadditivity} proves that it cannot merge two independent variables additively.  Theorem~\ref{thm:scalar-characterization} then gives the exact scalar fragment.

This also explains why the earlier analytic candidates felt over-engineered.  They were trying to hide two sorts inside one scalar formula.  Once scalar terms and pair terms are separated, the roles become simple:
\begin{center}
\begin{tabular}{p{0.24\linewidth}p{0.34\linewidth}p{0.32\linewidth}}
\toprule
Layer & Native behavior & Representative construction \\
\midrule
Visible scalar core & predecessor-scaling, contraction, multiplication & $\B(x,y)=xy-y$ \\
Pair/shadow layer & mixed addition & $\Tr\langle x\mid y\rangle=x+y$ \\
Derived projections & norm and spread & $\Nm,\Sp$ as macros \\
\bottomrule
\end{tabular}
\end{center}

Finally, the viewpoint suggests a useful diagnostic for symbolic search.  If a scalar grammar persistently fails to generate a target, the right response may not be to complicate the scalar operator.  The failure may indicate that the target lives in another sort or appears only after a projection.

\section{Beyond the golden ratio: quadratic units}\label{sec:quadratic-units}

The golden construction is a special case of a broader quadratic-unit pattern.  Let $\alpha$ be a root of
\[
    t^2-kt-d=0,
\]
with $k,d\in\Z$.  We emphasize the unit cases $d=\pm1$ because
\[
    B_k(\alpha,\alpha)=d,
\]
so $d=1$ gives the rational unit directly in the scalar layer, while $d=-1$ gives its negative and hence the unit after scalar negation.  The formal template works for any $d$, but the unit cases are the closest analogues of the golden seed.  Its conjugate is
\[
    \alpha^\sig=k-\alpha,
\]
and the defining invariants are
\[
    \Tr(\alpha)=k,
    \qquad
    \Nm(\alpha)=-d,
    \qquad
    (\alpha-\alpha^\sig)^2=k^2+4d.
\]
Under a chosen embedding, the last relation fixes a corresponding square-root sign.

The scalar core attached to the trace parameter $k$ is
\[
    B_k(x,y)=y(x-k).
\]
Then the defining polynomial immediately gives
\[
    B_k(\alpha,\alpha)=\alpha(\alpha-k)=d.
\]

\begin{proposition}[Quadratic-unit template]\label{prop:quadratic-template}
Let $\alpha$ satisfy $\alpha^2-k\alpha-d=0$ and let $\alpha^\sig=k-\alpha$.  With the pair
\[
    A=\langle \alpha\mid \alpha^\sig\rangle,
\]
the shifted scalar core and shadow projections satisfy
\[
    B_k(\alpha,\alpha)=d,
    \qquad
    \Tr(A)=k,
    \qquad
    \Nm(A)=-d,
\]
and the derived spread is
\[
    \Sp(A)=\alpha-\alpha^\sig=2\alpha-k,
\]
so
\[
    \Sp(A)^2=k^2+4d.
\]
Under a chosen embedding, this may be written with the corresponding square-root sign.
\end{proposition}

This template makes the golden, silver, and Gaussian cases variations of the same construction.

\begin{center}
\begin{tabular}{llllll}
\toprule
Field & Leaf $\alpha$ & Polynomial & $(k,d)$ & Scalar core & $(\Tr,\Nm,\Sp)$ \\
\midrule
Golden & $\ph$ & $t^2-t-1$ & $(1,1)$ & $y(x-1)$ & $(1,-1,\sqrt5)$ \\
Silver & $1+\sqrt2$ & $t^2-2t-1$ & $(2,1)$ & $y(x-2)$ & $(2,-1,2\sqrt2)$ \\
Gaussian & $i$ & $t^2+1$ & $(0,-1)$ & $xy$ & $(0,1,2i)$ \\
\bottomrule
\end{tabular}
\end{center}

For the silver ratio $\delta=1+\sqrt2$, the conjugate is $\delta^\sig=1-\sqrt2=-\delta^{-1}$ and $B_2(\delta,\delta)=1$.  The shadow pair $\Delta=\langle\delta\mid\delta^\sig\rangle$ gives
\[
    \Tr(\Delta)=2,
    \qquad
    \Nm(\Delta)=-1,
    \qquad
    \Sp(\Delta)=2\sqrt2.
\]
For the Gaussian integer $i$, the minimal polynomial is $t^2+1=t^2-0t-(-1)$, so $k=0$ and $d=-1$.  The scalar core collapses to ordinary multiplication:
\[
    B_0(x,y)=xy.
\]
Here $B_0(i,i)=-1$, while the shadow pair $I=\langle i\mid -i\rangle$ gives
\[
    \Tr(I)=0,
    \qquad
    \Nm(I)=1,
    \qquad
    \Sp(I)=2i.
\]

\section{Open questions}

\begin{question}[General scalar characterizations]
Theorem~\ref{thm:scalar-characterization} resolves the two-variable golden scalar fragment: it is the multiplicative, constant-shift closure of $R[x]$ and $R[y]$, where $R=\Z[\ph]$.  Does the same characterization extend cleanly to more variables, and to shifted scalar cores
\[
    B_k(f,g)=g(f-k)
\]
over other quadratic-unit rings?
\end{question}

\begin{question}[Minimality of trace completion]
Theorem~\ref{thm:completion} shows that pair formation and $\Tr$ complete the ordinary ring operations.  Is this completion minimal in a precise many-sorted sense?  For example, can one prove that no unary scalar extension, or no pair operation weaker than trace, closes the scalar non-additivity gap?
\end{question}

\begin{question}[Term equivalence]
Is the typed shadow algebra term-equivalent, in an appropriate many-sorted sense, to a standard presentation of commutative ring arithmetic over $\Z[\ph]$ or $\Q(\ph)$?
\end{question}

\begin{question}[Higher-degree shadow layers]
For a degree-$n$ number field, should the shadow layer use an $n$-slot object
\[
    \langle x_1\mid x_2\mid\cdots\mid x_n\rangle?
\]
If the goal is only to complete ordinary ring arithmetic, a trace projection may still be the only primitive additive channel: products can be built in the scalar layer, and elementary symmetric functions can be formed by tracing tuples of squarefree products.  For example, in a cubic shadow layer,
\[
    e_1=x+y+z=\Tr\langle x\mid y\mid z\rangle,
\]
while
\[
    e_2=xy+xz+yz=\Tr\langle xy\mid xz\mid yz\rangle,
\]
and
\[
    e_3=xyz
\]
is scalar-multiplicative.  If the goal is instead to reconstruct minimal polynomials or resolvents from opaque conjugate tuples, then the full elementary-symmetric projection tower $(e_1,e_2,\dots,e_n)$ may be the natural primitive structure.
\end{question}

\begin{question}[Analytic extension]
Can a simple branch or analytic layer be added to the shadow algebra to recover transcendental functions, in analogy with EML?  Or is the construction best understood as algebraic rather than analytic?
\end{question}

\section{Conclusion}

EML begins with the rational unit $1$ and shows that a single analytic binary operator can generate the elementary functions.  The present note explores a different seed: the golden ratio $\ph$, an algebraic unit whose defining identity already contains the rational unit:
\[
    \ph^2-\ph=1.
\]

A clean scalar operation emerges:
\[
    \B(x,y)=xy-y.
\]
With leaf $\ph$, this operation generates $1$, $0$, $\ph^{-1}$, $\ps$, negation, predecessor, successor, and multiplication.  The scalar fragment is now characterized exactly: it consists of multiplicative trees built from one-variable polynomials, with constant shifts from $\Z[\ph]$.  In particular, scalar addition is provably absent from the grammar; Theorem~\ref{thm:scalar-nonadditivity} rules out mixed affine-linear terms such as $x+y$, and the descent criterion rules out many higher-degree separable sums as well.  Rather than treating this as a failure, we interpret it as evidence that addition is not native to the visible scalar layer.

The scalar constants are exactly $\Z[\ph]$.  This ring is countable, so it cannot be the real line.  But because $\ph$ is irrational, $\Z[\ph]$ is dense in $\R$ under the chosen real embedding, while remaining discrete under the full Minkowski embedding.  Thus the finite grammar does not generate every real number exactly; instead, in one real coordinate it generates a dense golden arithmetic whose metric completion is $\R$.

The conjugate shadow supplies the missing layer.  Pair terms $\langle x\mid y\rangle$ together with the trace projection complete the ring operations:
\[
    x+y=\Tr\langle x\mid y\rangle,
    \qquad
    xy=\B(x+1,y).
\]
Norm and spread are then derived macros.  For the golden pair $\Phi=\langle\ph\mid\ps\rangle$, these projections recover
\[
    \Tr(\Phi)=1,
    \qquad
    \Nm(\Phi)=-1,
    \qquad
    \Sp(\Phi)=\sqrt5.
\]

The same architecture extends to quadratic units through the shifted scalar cores $B_k(x,y)=y(x-k)$, with the golden, silver, and Gaussian cases appearing as variations of the same trace/norm template.  Thus the proposed golden shadow algebra is not a replacement for EML's one-sorted universality theorem.  It is a complementary algebraic picture: a $\ph$-leaf system in which multiplication and predecessor-scaling live in the visible scalar core, addition appears naturally as trace from a conjugate-pair shadow, and the real line appears only after completion.

\begin{thebibliography}{9}

\bibitem{Odrzywolek2026EML}
A.~Odrzywo\l{}ek.
\newblock \emph{All elementary functions from a single binary operator}.
\newblock arXiv:2603.21852, 2026.

\bibitem{Hurwitz1891}
A.~Hurwitz.
\newblock Ueber die angenäherte Darstellung der Irrationalzahlen durch rationale Brüche.
\newblock \emph{Mathematische Annalen}, 39:279--284, 1891.

\bibitem{HardyWright2008}
G.~H. Hardy and E.~M. Wright.
\newblock \emph{An Introduction to the Theory of Numbers}.
\newblock 6th edition, Oxford University Press, 2008.

\bibitem{Marcus1977}
D.~A. Marcus.
\newblock \emph{Number Fields}.
\newblock Springer-Verlag, 1977.

\bibitem{ConradTraceNorm}
K.~Conrad.
\newblock \emph{Trace and Norm}.
\newblock Expository notes, University of Connecticut. Available at \url{https://kconrad.math.uconn.edu/blurbs/galoistheory/tracenorm.pdf}.

\bibitem{Lawvere1963}
F.~W. Lawvere.
\newblock \emph{Functorial Semantics of Algebraic Theories}.
\newblock Ph.D. thesis, Columbia University, 1963.

\bibitem{HylandPower2007}
M.~Hyland and J.~Power.
\newblock The category theoretic understanding of universal algebra: Lawvere theories and monads.
\newblock \emph{Electronic Notes in Theoretical Computer Science}, 172:437--458, 2007.

\end{thebibliography}

\end{document}
